\section{The System as Built}
\label{sec:architecture}

\noindent\textbf{The rig.}
The vehicle side is an ordinary Android phone and a Jetson Orin Nano, with no
purpose-built hardware (Figure~\ref{fig:overview}(b)). The phone holds no
interpretation logic. It captures or queries at the rate it is told and forwards
raw data, and every association, fusion and sampling decision happens on the
Jetson. The phone is connected to the internet via 5G cellular connection.
Two transports are implemented for the same wire protocol. A Tailscale network
is used for development with the phone in hand, and USB is used in the vehicle.
Both devices are powered from the vehicle, so energy is not a binding
constraint.

\noindent\textbf{The sensing scheduler.}
Each sensor's sampling rate is set at run time by the Jetson and commanded to
the phone. The four rates are camera, GPS, inertial unit and traffic feed. The
costs the scheduler trades against are traffic-API quota, thermal headroom on
the phone, and Jetson compute. It has two modes. In shadow mode the decision
is recorded and not sent, and in live mode the phone applies it. That
distinction is about sensing configuration only. The rate command carries
sampling rates and a shadow flag. The
word ``live'' here describes what the sensing configuration
did, and not what the vehicle did.

\noindent\textbf{The transport.}
One wire protocol carries everything between the phone and the Jetson, over the
development link and the in-vehicle link alike. It defines eight channels, each
with its own sequence numbers, priority, overflow policy and queue depth, over a
shared timebase. The frame layout is frozen in a single specification file.
Neither implementation is tested against the other. Both are tested against that
file, so the two sides cannot drift by agreeing with each other.

\noindent\textbf{The safety gate.}
The advisory passes a filter before it is displayed. The filter is vendored
onto the device from the simulator's safety layer and checked against a frozen
copy of it, so the two cannot drift apart. The simulator does not run the layer
itself, and Section~\ref{sec:simulation} gives the reason. The advisory is a
speed, so a rule earns its place by bounding that speed, and two do. The first
raises a recommendation that would have the vehicle travel slowly on an
uncongested road, which is the obstruction case. The second overrides the
recommendation when the gap to the vehicle ahead, or the time left to close it,
falls below the configured minimum. Each rule records, per tick, whether it
fired, stayed quiet, or could not be evaluated at all, with the missing inputs
named. A rule that cannot be evaluated is inert, so it cannot move the
recommendation in either direction.

\noindent\textbf{The device cannot drift from what was trained.}
The observation and action contract is vendored onto the device from the
simulator and pinned to a named version of it, so the device encodes an
observation and decodes an action exactly as training did. The export step
refuses any dimension mismatch. We enforce this using a second frozen specification file, holding
encoded observations, the action schema, the decoders and the neutral
fallbacks. Every entry in it is derived twice, once from the pinned version of
the simulator and once from the vendored copy, and the generator refuses to
write the file on any disagreement. That derivation runs where the training
stack lives. The check that runs on the device is the comparison against the
frozen file, which needs the file and nothing else, so the property is
testable on the hardware that has to hold it.

\noindent\textbf{Two design stances.}
The first is to degrade rather than stop. Missing GPS, no camera, no display and
no trained checkpoint each have a defined degraded mode, because debugging time
inside a vehicle is expensive. The second is that the hot path holds no queues,
and the latest value wins. A frame that was never processed is dropped, a slow
display skips ticks, and a stalled card drops log records. For a real-time
advisory, stale data is worse than missing data.

\noindent\textbf{Data collection.} A tick is one pass of the Jetson's pipeline. 
It takes the most recent camera
frame together with the latest GPS and inertial readings, runs detection,
tracking and distance estimation, builds the observation, runs the policy,
decodes the advisory, passes it through the safety gate, and writes one record.
Camera frames pace the ticks, so the tick rate follows the camera rate the
sensing scheduler has commanded. Every per-tick count in this paper refers to
that record.